%%
%% This is file `hep-text-documentation.tex',
%% generated with the docstrip utility.
%%
%% The original source files were:
%%
%% hep-text-implementation.dtx  (with options: `documentation')
%% This is a generated file.
%% Copyright (C) 2019-2023 by Jan Hajer
%% This file may be distributed and/or modified under the
%% conditions of the LaTeX Project Public License, either
%% version 1.3c of this license or (at your option) any later
%% version. The latest version of this license is in:
%% http://www.latex-project.org/lppl.txt
%% and version 1.3c or later is part of all distributions of
%% LaTeX version 2005/12/01 or later.

\ProvidesFile{hep-text-documentation.tex}[2024/11/01 v1.3 HEP-Text documentation]

\RequirePackage[l2tabu, orthodox]{nag}
\documentclass{ltxdoc}

\renewcommand\theCodelineNo{\rmfamily\tstyle\footnotesize\arabic{CodelineNo}}
\AtBeginEnvironment{macrocode}{\renewcommand{\ttdefault}{clmt}}
\renewcommand{\MacroFont}{\codestyle}
\AtBeginDocument{\DeleteShortVerb{\|}}
\AtBeginDocument{\MakeShortVerb{\"}}
\EnableCrossrefs
\CodelineIndex
\RecordChanges

\usepackage{hologo}

\usepackage[parskip,oldstyle,font=10pt]{hep-paper}
\bibliography{bibliography}

\acronym{CM}{computer modern}
\acronym{LM}{latin modern}

\newenvironment{columns}[1][.5]{%
  \par\vspace{-\bigskipamount}%
  \begin{minipage}[t]{\linewidth}%
  \begin{minipage}[t]{#1\linewidth}%
  \def\column{%
    \end{minipage}%
    \begin{minipage}[t]{\linewidth-#1\linewidth}%
  }%
}{\end{minipage}\end{minipage}\par}

\GetFileInfo{hep-text.sty}

\title{The \software{hep-text} package\thanks{This document corresponds to \software{hep-text}~\fileversion.}}
\subtitle{List and text extensions}
\author{Jan Hajer \email{jan.hajer@tecnico.ulisboa.pt}}
\date{\filedate}

\begin{document}

\newgeometry{vscale=.8, vmarginratio=3:4, includeheadfoot, left=11em, marginparwidth=4.6cm, marginparsep=3mm, right=7em}

\maketitle

\begin{abstract}
The \software{hep-text} package extends \hologo{LaTeX} lists using the \software{enumitem} package and provides some text macros.
\end{abstract}

The package can be loaded by "\usepackage{hep-text}".

\DescribeMacro{lang}
The "lang" option sets the used language and takes the values allowed by the \software{babel} package \cite{babel}, that is loaded for its hyphenation support.

Quotation commands are provided by the \software{csquotes} package \cite{csquotes}.
\DescribeMacro{\enquote}
\DescribeMacro{\MakeOuterQuote}
It provides the convenient macros "\enquote"\marg{text} and "\MakeOuterQuote{}" allowing to leave the choice of quotation marks to \hologo{LaTeX} and use \verb|"| instead of the pair \verb|``| and \verb|''|, respectively.

\DescribeMacro{\eg}
\DescribeMacro{\vs}
The \software{foreign} package \cite{foreign} defines macros such as "\eg", "\ie", "\cf", and "\vs" which are typeset as \eg, \ie, \cf, and \vs with the appropriate spacing.
Issuing "\renewcommand\foreignabbrfont{\itshape}" these abbreviations are typeset in italic.

\DescribeMacro{\no}
The "\no"\marg{number} macro is typeset as \no{123}.

\DescribeMacro{\software}
The "\software"\oarg{version}\marg{name} macro is typeset as \software[\fileversion]{HEP-Paper}.

\DescribeMacro{\online}
The "\online"\marg{url}\marg{text} macro combines the features of the "\href"\marg{url}\allowbreak\marg{text}\allowbreak \cite{hyperref} and the "\url"\marg{text} \cite{url} macros, resulting in \eg \online{https://ctan.org/pkg/hep-text}{ctan.org/pkg/hep-text}.

\DescribeMacro{inlinelist}
\DescribeMacro{enumdescript}
The "inlinelist" and "enumdescript" environments are defined using the \software{enumitem} package \cite{enumitem}.
\\
\begin{columns}
\begin{verbatim}
The three main points are
\begin{inlinelist}
  \item one
  \item two
  \item three
\end{inlinelist}
\end{verbatim}
\column
The three main points are
\begin{inlinelist}
 \item one
 \item two
 \item three
\end{inlinelist}
\end{columns}
\vspace{4ex}
\begin{columns}[.6]
\begin{verbatim}
\begin{enumdescript}[label=\Roman*)]
  \item{First} one
  \item{Second} two
  \item{Third} three
\end{enumdescript}
\end{verbatim}
\column
\begin{enumdescript}[label=\Roman*)]
 \item{First} one
 \item{Second} two
 \item{Third} three
\end{enumdescript}
\end{columns}

\DescribeMacro{\underline}
The "\underline" macro is redefined to allow line-breaks using the \software{soul} package \cite{soul}.

\printbibliography

\end{document}

\endinput
%%
%% End of file `hep-text-documentation.tex'.
